\documentclass[11pt]{article}
\usepackage[margin=1in]{geometry}
\usepackage{amsmath,amssymb,amsthm}
\usepackage{hyperref}
\hypersetup{hidelinks}

\newtheorem{theorem}{Theorem}
\newtheorem{lemma}[theorem]{Lemma}
\newtheorem{corollary}[theorem]{Corollary}
\theoremstyle{definition}
\newtheorem{definition}{Definition}

\setlength{\parskip}{4pt plus 1pt minus 1pt}
\setlength{\parindent}{0pt}

\title{Evaluation-Driven State Revision and the\\Encoding Hierarchy Under Bounded Context}
\author{Patrick D.\ McCarthy}
\date{}

\begin{document}
\maketitle

\begin{abstract}
We prove two linked results for any entity persisting under bounded
active context with survival pressure. First, agency ($\gamma > 0$)
necessarily produces evaluation-driven state revision --- gradient
descent on uncertainty $U$ --- with convergence guaranteed by a scalar
Lyapunov argument: $U \geq 0$ is bounded below, each update reduces
$\mathbb{E}[U]$ via the data processing inequality, and monotone
convergence applies without dimensional decomposition. The gradient
never terminates because the reachable trajectory space grows
monotonically. Second, the certainty threshold for encoding an action
at level $L_i$ is $\theta_i = 1 - \varepsilon / C_i$, where $C_i$ is
the cost of failure at that level. This derives the learning rate
hierarchy: evolutionary rates are necessarily orders of magnitude lower
than individual learning rates, not by assumption but from
cost-of-failure differences.
\end{abstract}

% ------------------------------------------------------------------
\section{Introduction}
\label{sec:intro}
% ------------------------------------------------------------------

An entity with agency $\gamma > 0$ acts on the world and receives feedback.
What is the necessary structure of learning? Papers~1--3 in this
series established that any entity persisting under bounded active
context must maintain a directed graph over propositions
\cite{McCarthy2026a}, that escalation dominates top-down allocation
\cite{McCarthy2026b}, and that knowledge and action are inseparable
under survival pressure \cite{McCarthy2026c}. What remains open is two
questions: what \emph{dynamics} govern the updating process, and what
determines the \emph{threshold} at which an action earns promotion to
a more permanent encoding level.

This paper answers both. We prove that the updating process is
gradient descent on uncertainty $U$, with convergence guaranteed by a
scalar Lyapunov argument that requires no dimensional decomposition
(Theorem~\ref{thm:gradient}). We then prove that the certainty
threshold for encoding at level $L_i$ is determined by the cost of
failure at that level, which derives the learning rate hierarchy across
encoding levels (Theorem~\ref{thm:encoding}).

Several lines of prior work converge on these questions from different
angles. Dupoux, LeCun, and Malik \cite{DupouxLeCunMalik2026} propose
a bilevel optimization framework (evolution as outer loop, development
as inner loop) but treat the learning rate relationship between levels
as a design parameter rather than deriving it. Cover and Thomas
\cite{Cover2006} provide the data processing inequality that anchors
our convergence argument. Schmidhuber \cite{Schmidhuber2010} identifies
curiosity as compression progress --- structurally identical to our
learning rate $\eta_M$. Hutter \cite{Hutter2005} constructs AIXI as
an idealized agent but assumes unbounded computation, precisely the
constraint our framework takes as central. Chollet
\cite{Chollet2019} defines intelligence as skill-acquisition
efficiency --- close to $\eta_M$ but treating prior knowledge as
confound rather than constitutive element. Hinton and Nowlan
\cite{Hinton1987} demonstrate that learning can guide evolution (the
Baldwin effect) without formalizing the threshold that governs
promotion. Finn et al.\ \cite{FinnEtAl2017} and Thrun
\cite{Thrun1998} develop meta-learning algorithms that operate as $M$
on $M$ --- the $M^{(2)}$ level in our hierarchy --- but do not derive
why the meta-level necessarily learns more slowly.

% ------------------------------------------------------------------
\section{Formal Framework}
\label{sec:framework}
% ------------------------------------------------------------------

We introduce the minimal definitions required for the two main
results, extending the framework of \cite{McCarthy2026a}.

\begin{definition}[Entity, World, Uncertainty]
\label{def:entity}
An entity $E = (K, A, \sigma, \gamma)$ consists of a knowledge base $K$
(propositions about world $W$), an action space $A$ (executable
actions), a selection mechanism $\sigma$ that chooses actions given
context, and an agency parameter $\gamma$. The world state space $W$ is a
measurable space; the true state $w \in W$ is not directly observable.
The \emph{uncertainty} of $E$ about state $w$ given knowledge $K$ is:
\[
U(w, K) = -\log P(w \mid K)
\]
The expected uncertainty is the conditional entropy:
\[
H(W \mid K) = \mathbb{E}_w[U(w, K)] = -\sum_w P(w \mid K) \log P(w \mid K)
\]
\end{definition}

\begin{definition}[Agency]
\label{def:agency}
The agency parameter $\gamma \in [0,1]$ is the probability that $E$
engages its learning mechanism $M$ in response to a gradient signal.
At $\gamma = 0$, $A$ and $K$ are fixed forever: the entity receives
signals but never updates. Operationally:
\[
\gamma(E) \;\propto\; \lim_{U_{\mathrm{lethal}} \to \infty}
\frac{dU}{dt}
\]
That is, agency measures whether the entity continues engaging $M$
when survival pressure is removed. An entity that updates only under
immediate threat has $\gamma \approx 0$; an entity that updates
spontaneously has $\gamma > 0$.
\end{definition}

\begin{definition}[Action Feedback]
\label{def:feedback}
For action $a$ executed in world state $w$, the \emph{feedback}
$\varphi(a, w)$ is the observable consequence. The entity updates its
knowledge base: $K_a = K \cup \{p_a\}$, where $p_a$ is the
proposition encoding the feedback. The \emph{gradient signal} is the
prediction error:
\[
\delta = \varphi(a, w) - \mathbb{E}[\varphi(a, \cdot) \mid K]
\]
\end{definition}

\begin{definition}[Information Gain]
\label{def:gain}
The information gain from executing $a$ given $K$ is the mutual
information between the action's outcome and the world state,
conditioned on current knowledge:
\[
G(a, K) = I(a; W \mid K) \geq 0
\]
The non-negativity follows from the data processing inequality
\cite{Cover2006}. $G = 0$ if and only if $a$ provides no information
about $W$ beyond what $K$ already contains.
\end{definition}

\begin{definition}[Encoding Hierarchy]
\label{def:hierarchy}
The encoding space $\mathcal{E}(c, r)$ is a continuous space
parameterized by cost $c$ and reversibility $r$. Within this space,
interfaces $L_0, L_1, \ldots, L_n$ define discrete encoding levels
with increasing cost and reversibility:
\begin{itemize}
\item $L_0$ (genomic): highest cost of error, lowest reversibility,
   slowest update rate.
\item $L_n$ (active context): lowest cost of error, highest
   reversibility, fastest update rate.
\end{itemize}
Escalation dominance \cite{McCarthy2026b} pushes proven actions
upward toward more permanent levels. The present paper derives the
threshold that governs this promotion.
\end{definition}

% ------------------------------------------------------------------
\section{Gradient Descent on Uncertainty}
\label{sec:gradient}
% ------------------------------------------------------------------

\begin{theorem}[Evaluation-Driven State Revision]
\label{thm:gradient}
Let $E = (K, A, \sigma, \gamma)$ be an entity with $\gamma > 0$ persisting in
world $W$ under bounded active context $C_n$.

\textbf{Part I.} Every executed action generates a gradient signal,
and $\gamma > 0$ guarantees engagement with that signal, producing
evaluation-driven state revision --- gradient descent on $U$ --- with
convergence to a local minimum.

\textbf{Part II.} The gradient is never globally zero: the reachable
trajectory space grows monotonically with $K$, so descent never
terminates within the reachable epistemic domain.
\end{theorem}

\begin{proof}
\textbf{Part I: Forced Signal, Response Gate, and Convergence.}

Every executed action $a$ generates feedback $\varphi(a, w)$ from the
world. This signal arrives whether or not the entity has $\gamma > 0$: the
world responds to actions regardless of the actor's capacity to learn
from them. The gradient signal $\delta = \varphi(a, w) -
\mathbb{E}[\varphi(a, \cdot) \mid K]$ is therefore produced at every
action.

Agency $\gamma$ is the response gate. An entity with $\gamma = 0$ receives
every signal $\delta$ but never updates: $(A, K)$ remain static. An
entity with $\gamma > 0$ engages $M$ with probability $\gamma$ on each signal,
producing the update:
\[
\Delta_a = -\eta \cdot \delta
\]
where $\eta > 0$ is the learning rate of mechanism $M$. The update
$\Delta_a = -\eta \cdot \delta$ is a stochastic gradient step in the
sense of Robbins and Monro \cite{RobbinsMonro1951}: $\delta$ is a
noisy estimator of $\nabla_K U$, with the noise arising from the
stochasticity of $w$. The data processing inequality (Step~1 below)
guarantees that $\mathbb{E}[\Delta_a]$ is in a descent direction,
which is the defining property of stochastic gradient descent.

We now establish convergence via a scalar Lyapunov argument that
requires no dimensional decomposition.

\emph{Step 1: Each $M$-update with $G > 0$ strictly reduces
$\mathbb{E}[U]$.} When $M$ integrates the feedback from $a$ into $K$,
forming $K' = K \cup \{p_a\}$, the data processing inequality
\cite{Cover2006} guarantees:
\[
H(W \mid K') \leq H(W \mid K)
\]
with equality if and only if $G(a, K) = 0$. When $G(a, K) > 0$:
\[
\mathbb{E}[U(w, K') \mid K] \leq U(w, K) - G(a, K)
\]
Each informative update strictly reduces expected uncertainty.

\emph{Step 2: Monotone convergence.} Consider the sequence
$\{E_t\} = \{\mathbb{E}[U(w, K(t))]\}$ indexed by update steps $t$.
By Step~1, $\{E_t\}$ is non-increasing. By definition, $U(w, K) =
-\log P(w \mid K) \geq 0$, so $\{E_t\}$ is bounded below by $0$. By
the monotone convergence theorem, $\{E_t\}$ converges. Since the
sequence is non-increasing and convergent:
\[
\sum_{t=0}^{\infty} G(a_t, K(t)) \leq U(w, K(0)) < \infty
\]
Therefore $G(a_t, K(t)) \to 0$ as $t \to \infty$. The entity
converges to a local minimum of $U$ --- a state where available
actions yield diminishing informational returns.

\emph{Step 3: Dimensionality irrelevance.} The entire argument
operates on the scalar quantity $U$. The entity's internal
parameterization may grow as the reachable trajectory space
$T_{\mathrm{reachable}}$ expands. Each new proposition $p_a$ added to
$K$ can only reduce $U$ (by the data processing inequality): knowledge
never increases expected uncertainty. Cross-grounding between
propositions accelerates convergence by allowing a single observation
to update multiple related beliefs. Adding structure to $K$ is
analogous to adding a column to a design matrix and re-normalizing:
the residual can only shrink. The Lyapunov bound $U \geq 0$ holds
regardless of the dimensionality of $K$ or the parameterization of
$M$.

\textbf{Part II: Non-Termination.}

The reachable trajectory space $T_{\mathrm{reachable}}(K)$ grows
monotonically with $K$. Executing any action $a$ produces feedback
that grounds a new proposition $p_a$, yielding $K_a = K \cup \{p_a\}$
with:
\[
T_{\mathrm{reachable}}(K_a) \supseteq T_{\mathrm{reachable}}(K)
\]
The inclusion is strict whenever $p_a$ enables new actions or reveals
new state dimensions. At any finite time $t$, the entity has explored
a finite subset of $T_{\mathrm{reachable}}(K(t))$. The uncovered
region always contains states $w'$ with $U(w', K(t)) > 0$. Therefore
there always exists some action $a'$ with $G(a', K(t)) > 0$.

The gradient is never globally zero. Descent converges locally (Part
I) but never terminates globally: the epistemic frontier recedes as
knowledge grows.

\emph{Interaction with the curation regime.} The convergence argument
in Part~I assumes $K$ grows monotonically ($K' = K \cup \{p_a\}$). In
the curation regime established in \cite{McCarthy2026c}, $|K|$ is
bounded at $K_{\max}$ and every addition requires displacement. Does
displacement break monotonicity? Not under the displacement criterion
of the $K/A$ inseparability theorem: proposition $p$ is displaced by
$q$ only when $\mathbb{E}[G(a_q, K)] > \mathbb{E}[G(a_p, K)]$ ---
strictly higher marginal value. The entity replaces less-grounding
propositions with more-grounding ones, preserving the descent
direction. Monotonicity of $H(W \mid K)$ reduction is maintained not
because $K$ grows, but because each displacement strictly improves the
information content of a fixed-size $K$.

\textbf{Remark.} An entity with $\gamma = 0$ never updates. As $W$
changes, its fixed $K$ becomes increasingly inaccurate: $U(w_t, K)$
grows until it exceeds the lethality threshold $U_{\mathrm{lethal}}$.
Selection removes such entities. Therefore $\gamma > 0$ is necessary for
persistence --- a claim established fully in Paper~5.
\end{proof}

% ------------------------------------------------------------------
\section{Multi-Timescale Structure}
\label{sec:multiscale}
% ------------------------------------------------------------------

\begin{corollary}[Individual Learning and Evolution Share Update Structure]
\label{cor:multiscale}
Both individual learning and evolutionary selection instantiate the
same abstract update: execute, receive signal, retain or revise. They
differ in mechanism, object, and timescale.
\end{corollary}

\begin{proof}
\textbf{Individual learning.} The entity updates its action
space directly from continuous gradient signals:
\[
\Delta a_{\mathrm{ind}} = -\eta_{\mathrm{ind}} \cdot \nabla_A U(w, K)
\]
The update is directed: $\delta$ specifies both magnitude and
direction of revision. The learning rate $\eta_{\mathrm{ind}}$ is
determined by the capacity of mechanism $M$.

\textbf{Evolutionary selection.} The population-level update operates
on a different object (the heritable component of $A$) via a different
mechanism:
\[
\Delta a_{\mathrm{species}} = \mathrm{select}(\mathrm{mutate}(A_{\mathrm{population}}))
\]
Mutation is undirected: random perturbation of heritable parameters.
Selection is binary: survive or not. No individual organism descends a
gradient at the evolutionary timescale --- the population distribution
shifts toward lower $U$ through differential survival.

Both processes move the entity (individual or lineage) toward lower
$U$ at their respective timescales. Gradient descent is the efficient
limit: directed $\delta$ combined with capable $M$ produces the
steepest feasible descent. Evolution is the general case: random
perturbation combined with selection produces descent without requiring
any individual to compute a gradient.

The two are not independent. Individual learning discovers actions at
$L_n$ that, if sufficiently validated, promote to $L_0$ via the
encoding hierarchy (Theorem~\ref{thm:encoding}). This is the Baldwin
effect \cite{Hinton1987}, here derived rather than observed.
\end{proof}

% ------------------------------------------------------------------
\section{Encoding Permanence}
\label{sec:encoding}
% ------------------------------------------------------------------

\begin{theorem}[Encoding Permanence Scales with Rigor]
\label{thm:encoding}
Let $L_0, \ldots, L_n$ be encoding levels with associated failure
costs $C_0 > C_1 > \cdots > C_n$ and reversibility costs
$r_0 > r_1 > \cdots > r_n$. Let $\varepsilon$ be the maximum
acceptable expected error cost. Then the minimum certainty threshold
for encoding an action at level $L_i$ is:
\[
\theta_i = 1 - \frac{\varepsilon}{C_i}
\]
\end{theorem}

\begin{proof}
Encoding action $a$ at level $L_i$ is effectively irreversible at cost
$r_i$. If the encoding is erroneous --- $a$ is not in fact reliably
adaptive --- the cost of that error is $C_i$. If $\theta$ is the
certainty that $a$ is correct, then the expected error cost is:
\[
\mathbb{E}[\text{error cost}] = (1 - \theta) \cdot C_i
\]
Encoding is safe when the expected error cost does not exceed the
acceptable threshold:
\[
(1 - \theta) \cdot C_i \leq \varepsilon
\]
Rearranging:
\[
\theta \geq 1 - \frac{\varepsilon}{C_i}
\]
The minimum safe threshold is therefore $\theta_i = 1 -
\varepsilon / C_i$.

\textbf{Consequence.} At $L_0$ (genomic encoding), $C_0$ is the cost
of lineage extinction --- orders of magnitude larger than $\varepsilon$.
Therefore $\theta_0 \approx 1$: near-certainty is required. A bad
encoding at $L_0$ does not kill one individual; it kills the lineage.
At $L_n$ (active context), $C_n$ is small (cost of a single failed
action), so $\theta_n$ is meaningfully less than $1$: the entity can
afford to try uncertain actions. Rigor is proportional to permanence.
\end{proof}

\textbf{Error propagation through the hierarchy.} A natural objection
holds that errors compound across levels: if each level has error rate
$1 - \theta_i$, the cumulative error grows multiplicatively. This
objection is based on a naive cascading model. Under escalation
\cite{McCarthy2026b}, errors partition into two categories:
\emph{handled errors}, which are resolved at level $L_i$ and never
reach $L_{i+1}$, and \emph{unhandled exceptions}, which represent
genuinely irreducible failures that escalate. The error rate seen by
$L_{i+1}$ is not a compounded Bernoulli product of lower-level error
rates but the rate of genuinely irreducible errors --- failures that
persist across sufficient variation to pass the threshold $\theta_i$.
The thresholds $\theta_i$ are independent by construction: each is
determined by $C_i$ and $\varepsilon$, not by the error rates of
adjacent levels.

\subsection{Deriving the Learning Rate Hierarchy}

The learning rate at level $L_i$ is inversely related to the rigor
threshold. To see why, consider the sample complexity: encoding at
$L_i$ requires confidence $\theta_i$ that the action is adaptive. By
Hoeffding's inequality \cite{Hoeffding1963}, the number of independent
trials $n_i$ needed to achieve confidence $\theta_i$ satisfies:
\[
n_i \;\geq\; \frac{\log(1/(1 - \theta_i))}{2\varepsilon^2}
\;\approx\; \frac{1}{2\varepsilon^2} \log \frac{C_i}{\varepsilon}
\]
The effective learning rate --- the fraction of experience that results
in permanent encoding --- is inversely proportional to $n_i$:
\[
\eta_i \;\propto\; \frac{1}{n_i}
\;\propto\; \frac{1}{\log(C_i / \varepsilon)}
\]

For $L_n$ (active context): $C_n$ is small but finite,
$\theta_n = 1 - \varepsilon / C_n$ is meaningfully less than $1$, and
$\eta_{\mathrm{ind}}$ is large. The entity updates rapidly from each
gradient signal.

For $L_0$ (genomic): $C_0 \gg \varepsilon$, so $\theta_0 \approx 1$
and $\eta_{\mathrm{evo}} \approx \varepsilon / C_0 \approx 0$. The
lineage updates only when certainty is overwhelming.

Since $\theta_0 \gg \theta_n$, it follows that
$\eta_{\mathrm{evo}} \ll \eta_{\mathrm{ind}}$. The learning rate
difference between evolution and individual learning is \emph{derived}
from the cost-of-failure structure, not assumed as a design parameter.

\begin{corollary}[Bounded Adaptation]
\label{cor:bounded}
An action $a$ is \emph{certain} at level $L_i$ when $U(w, K_a)
\approx 0$ across sufficient variation to meet threshold $\theta_i$.
Certain actions escape active context permanently: they are encoded at
$L_i$ and no longer require $L_n$ resources to execute. This frees
active context capacity for the next epistemic frontier.
\end{corollary}

\begin{definition}[Certainty]
\label{def:certain}
Action $a$ is \emph{certain} when $U(w, K_a) \approx 0$ across
sufficient variation in $w$ to meet the threshold $\theta_i$ for
encoding level $L_i$. Formally: $P(\text{$a$ is adaptive} \mid K)
\geq \theta_i$, where the probability is taken over the distribution
of world states relevant to $L_i$.
\end{definition}

% ------------------------------------------------------------------
\section{Engagement with Prior Work}
\label{sec:prior}
% ------------------------------------------------------------------

\textbf{Dupoux, LeCun, and Malik \cite{DupouxLeCunMalik2026}.}
Their bilevel optimization framework --- inner loop as development,
outer loop as evolution --- captures the correct structural
relationship between encoding levels. However, they treat the learning
rate difference between levels as a design parameter: the inner loop
is fast because it is designed to be fast, and the outer loop is slow
because mutation and selection are slow. Theorem~\ref{thm:encoding}
derives this relationship from first principles: $\eta_{\mathrm{evo}}
\ll \eta_{\mathrm{ind}}$ because $C_0 \gg C_n$, not because of any
design choice.

Their framework includes a ``System~M'' that is ``not learned through
gradient descent'' but fixed by evolution. Our framework explains
\emph{why}: the rigor threshold $\theta_0 \approx 1$ means that any
component encoded at $L_0$ must be near-certain across all
environments in the species distribution. Only mechanisms that
generalize across the full environmental distribution meet this
threshold. System~M is not excluded from gradient descent by fiat; it
is excluded because the cost of a bad learning mechanism at $L_0$ is
lineage extinction, and the resulting threshold $\theta_0$ permits
only components validated across maximal variation.

\textbf{Hutter \cite{Hutter2005}: AIXI.}
AIXI constructs the theoretically optimal agent by enumerating all
computable hypotheses, weighted by Kolmogorov complexity. The
construction assumes unbounded computation --- precisely the
constraint our framework takes as central. Under bounded active
context $C_n$, the AIXI solution is unreachable: the entity cannot
enumerate hypotheses but must descend the gradient with finite
resources. Our Theorem~\ref{thm:gradient} characterizes what bounded
agents actually do: local gradient descent on $U$ with convergence to
local minima, not global optimization over hypothesis space.

\textbf{Chollet \cite{Chollet2019}: Intelligence as
skill-acquisition efficiency.}
Chollet's measure of intelligence --- the rate at which an agent
acquires new skills from limited experience --- is structurally close
to our learning rate $\eta_M$. The key divergence is in the treatment
of prior knowledge. Chollet treats priors as a confound to be
controlled for: intelligence is efficiency \emph{net of} prior
knowledge. In our framework, prior knowledge $K$ is not a confound
but a constitutive element: $\eta_M$ operates on $K$, and the
capacity of $M$ is inseparable from the structure of $K$
\cite{McCarthy2026c}. An entity with richer $K$ does not have
``unfair advantage'' --- it has done more gradient descent.

\textbf{Schmidhuber \cite{Schmidhuber2010}: Curiosity as compression
progress.}
Schmidhuber defines intrinsic motivation as the rate of compression
progress --- the reduction in description length of the agent's model
of the world. This is structurally identical to $\eta_M$: both
measure the rate of uncertainty reduction. The difference is
ontological. Schmidhuber treats curiosity as a design objective ---
something to be engineered into agents. Our framework suggests that
$\gamma > 0$ (the tendency to engage $M$ even without immediate survival
pressure) is not a design choice but a necessary condition for
persistence, a claim established fully in Paper~5.

\textbf{Hinton and Nowlan \cite{Hinton1987}: Learning guiding
evolution.}
Hinton and Nowlan demonstrate via simulation that individual learning
can accelerate evolutionary search --- the Baldwin effect. An
organism that can learn a trait during its lifetime effectively
smooths the fitness landscape for evolution. Our framework formalizes
this: individual learning at $L_n$ discovers actions via rapid
gradient descent ($\eta_{\mathrm{ind}}$ large). Encoding permanence
(Theorem~\ref{thm:encoding}) determines which of these actions
promote to $L_0$: only those meeting $\theta_0 \approx 1$ --- actions
validated across sufficient environmental variation to warrant
heritable encoding. The Baldwin effect is not an empirical curiosity
but a necessary consequence of the encoding hierarchy.

\textbf{Kirkpatrick et al.\ \cite{Kirkpatrick2017}: Catastrophic
forgetting and elastic weight consolidation.}
Kirkpatrick et al.\ demonstrate that neural networks trained
sequentially on multiple tasks catastrophically forget earlier tasks
unless protected by a consolidation mechanism. Their solution ---
elastic weight consolidation (EWC) --- assigns higher ``importance''
to weights encoding well-established knowledge, penalizing changes to
those weights during subsequent learning. This is precisely the
encoding hierarchy: weights at lower encoding levels (higher
$\theta_i$, higher cost of error) are protected; weights at higher
levels (lower $\theta_i$) are free to change. EWC's Fisher information
matrix serves as an empirical estimator of our cost-of-failure $C_i$.
The result provides direct empirical evidence that the encoding
hierarchy is not merely theoretical: artificial systems that violate
it (flat learning rate across all parameters) fail, and systems that
approximate it (EWC, progressive neural networks) succeed.

\textbf{Murray et al.\ \cite{Murray2014}: Cortical temporal
hierarchy.}
Murray et al.\ establish that intrinsic timescales of neural
activity increase systematically along the cortical hierarchy: sensory
cortex fluctuates rapidly, while prefrontal cortex maintains stable
representations over seconds. This gradient of temporal stability maps
directly onto the encoding hierarchy: fast-updating sensory areas
operate at $L_n$ (low cost of error, high reversibility), while
slow-updating prefrontal areas operate at lower encoding levels (high
cost of error, low reversibility). The empirical timescale gradient
matches the theoretical prediction $\eta_i \propto 1/\log(C_i /
\varepsilon)$: regions with higher cost of error update more slowly.

\textbf{Finn et al.\ \cite{FinnEtAl2017} and Thrun \cite{Thrun1998}:
Meta-learning.}
Model-Agnostic Meta-Learning (MAML) learns an initialization from
which task-specific learning converges rapidly. ``Learning to learn''
\cite{Thrun1998} frames this as optimizing the learning algorithm
itself. In our framework, meta-learning is $M$ operating on $M$ ---
the $M^{(2)}$ level. The meta-learner sits at a higher encoding level
than the base learner: it updates more slowly ($\eta_{M^{(2)}} <
\eta_M$) because the cost of a bad meta-learning rule ($C_{M^{(2)}}$)
exceeds the cost of a bad base-level update ($C_M$). The MAML inner
loop is $L_n$; the outer loop is $L_{n-1}$. The hierarchy extends:
evolution is meta-meta-$\ldots$-learning at $L_0$.

% ------------------------------------------------------------------
\section{Conclusion}
\label{sec:conclusion}
% ------------------------------------------------------------------

Gradient descent on uncertainty is not a metaphor for learning. It is
what survival-driven updating \emph{is}. Any entity with $\gamma > 0$
that persists under bounded context necessarily engages in
evaluation-driven state revision: executing actions, receiving
gradient signals, and updating knowledge to reduce $U$. The
convergence argument is scalar --- it operates on $U$ directly, with
no dependence on the dimensionality of the entity's internal
parameterization. The gradient never terminates because the reachable
trajectory space grows monotonically with knowledge.

The encoding hierarchy is not a design choice but a necessary
consequence of cost-of-failure differences across encoding levels.
The threshold $\theta_i = 1 - \varepsilon / C_i$ determines how
certain an action must be before it earns permanent encoding. From
this single equation, the entire learning rate hierarchy follows:
$\eta_{\mathrm{evo}} \ll \eta_{\mathrm{ind}}$ because
$C_0 \gg C_n$, not because evolution is ``inherently slow'' or
individual learning is ``inherently fast.''

Together with graph necessity \cite{McCarthy2026a}, escalation dominance
\cite{McCarthy2026b}, and K/A inseparability \cite{McCarthy2026c},
this paper completes the structural account of
how knowledge is acquired, validated, and encoded. What remains is to
establish that $\gamma > 0$ is itself necessary for persistence --- that
agency is not optional but required. This is the subject of Paper~5.

\bibliographystyle{plain}
\begin{thebibliography}{99}

\bibitem{Chollet2019}
F.~Chollet.
\newblock On the measure of intelligence.
\newblock \emph{arXiv preprint arXiv:1911.01547}, 2019.

\bibitem{Cover2006}
T.~M. Cover and J.~A. Thomas.
\newblock \emph{Elements of Information Theory}.
\newblock Wiley-Interscience, 2nd edition, 2006.

\bibitem{DupouxLeCunMalik2026}
E.~Dupoux, Y.~LeCun, and J.~Malik.
\newblock Why {AI} systems don't learn and what to do about it: Lessons on
  autonomous learning from cognitive science.
\newblock \emph{arXiv preprint arXiv:2603.15381}, 2026.

\bibitem{FinnEtAl2017}
C.~Finn, P.~Abbeel, and S.~Levine.
\newblock Model-agnostic meta-learning for fast adaptation of deep networks.
\newblock In \emph{Proceedings of the 34th International Conference on Machine
  Learning (ICML)}, pages 1126--1135, 2017.

\bibitem{Friston2010}
K.~Friston.
\newblock The free-energy principle: A unified brain theory?
\newblock \emph{Nature Reviews Neuroscience}, 11(2):127--138, 2010.

\bibitem{Hinton1987}
G.~E. Hinton and S.~J. Nowlan.
\newblock How learning can guide evolution.
\newblock \emph{Complex Systems}, 1(3):495--502, 1987.

\bibitem{Hoeffding1963}
W.~Hoeffding.
\newblock Probability inequalities for sums of bounded random variables.
\newblock \emph{Journal of the American Statistical Association},
  58(301):13--30, 1963.

\bibitem{Hutter2005}
M.~Hutter.
\newblock \emph{Universal Artificial Intelligence: Sequential Decisions Based
  on Algorithmic Probability}.
\newblock Springer, 2005.

\bibitem{Kirkpatrick2017}
J.~Kirkpatrick, R.~Pascanu, N.~Rabinowitz, J.~Veness, G.~Desjardins,
  A.~A. Rusu, K.~Milan, J.~Quan, T.~Ramalho, A.~Grabska-Barwinska,
  D.~Hassabis, C.~Summerfield, T.~Lillicrap, and D.~Kumaran.
\newblock Overcoming catastrophic forgetting in neural networks.
\newblock \emph{Proceedings of the National Academy of Sciences},
  114(13):3521--3526, 2017.

\bibitem{McCarthy2026a}
P.~D. McCarthy.
\newblock Graph structure is necessary for information preservation under
  bounded context.
\newblock Working paper, 2026.

\bibitem{McCarthy2026b}
P.~D. McCarthy.
\newblock Escalation dominates top-down allocation under bounded context.
\newblock Working paper, 2026.

\bibitem{McCarthy2026c}
P.~D. McCarthy.
\newblock Knowledge and action are inseparable under survival pressure.
\newblock Working paper, 2026.

\bibitem{Murray2014}
J.~D. Murray, A.~Bernacchia, D.~J. Freedman, R.~Romo, J.~D. Wallis,
  X.~Cai, C.~Padoa-Schioppa, T.~Pasternak, H.~Seo, D.~Lee, and
  X.-J. Wang.
\newblock A hierarchy of intrinsic timescales across primate cortex.
\newblock \emph{Nature Neuroscience}, 17(12):1661--1663, 2014.

\bibitem{RobbinsMonro1951}
H.~Robbins and S.~Monro.
\newblock A stochastic approximation method.
\newblock \emph{The Annals of Mathematical Statistics}, 22(3):400--407, 1951.

\bibitem{Schmidhuber2010}
J.~Schmidhuber.
\newblock Formal theory of creativity, fun, and intrinsic motivation
  (1990--2010).
\newblock \emph{IEEE Transactions on Autonomous Mental Development},
  2(3):230--247, 2010.

\bibitem{Shannon1948}
C.~E. Shannon.
\newblock A mathematical theory of communication.
\newblock \emph{Bell System Technical Journal}, 27(3):379--423, 1948.

\bibitem{Simon1955}
H.~A. Simon.
\newblock A behavioral model of rational choice.
\newblock \emph{The Quarterly Journal of Economics}, 69(1):99--118, 1955.

\bibitem{Thrun1998}
S.~Thrun and L.~Pratt.
\newblock \emph{Learning to Learn}.
\newblock Kluwer Academic Publishers, 1998.

\end{thebibliography}

\end{document}
